\begin{afterword}
\summaryhead

The post contrasts two ways of seeing the rules of rationality. Traditional Rationality states them as social obligations: whoever asks you to accept a belief owes you evidence, and a theory owes the world bold, falsifiable predictions. People catch social cheating more easily than breaches of abstract rules, so this phrasing spreads well. But it invites a strange defence. Religious believers say ``you can't justify your belief in science!'', as if showing the critic to be a hypocrite would excuse them.

On the Bayesian view, the rules are laws like the second law of thermodynamics. A reliable belief needs evidence as a refrigerator needs energy. A self-cooling refrigerator, or an accurate map of New York drawn without visiting, is possible only with negligible probability. Pointing at others who break the rule changes nothing: two badly designed engines both fail. The law may give everyone a right to their beliefs, but nature gives no one a right to accuracy, and no vote or authority can grant an exemption. There is only cause and effect.

\responsehead

The post's central point is correct and well put. That someone else also holds beliefs without evidence does not make yours accurate. This is the standard answer to ``you do it too'' (the tu quoque), and the post adds a reason for it: accuracy depends on causes, not on fairness, so another person's fault cannot help you. The engine example states this exactly.

The claim that these rules are laws like thermodynamics is made here by analogy, and the analogy fits some beliefs better than others. Its tiny probabilities come from detail: a blind guess at a map of New York must get every street right. A guess made without evidence about a single yes-or-no question is right about half the time. What it lacks is reliability, and the post's own statement of the law is about reliable belief. The post does not say how much detail the beliefs it has in mind carry, so the reader cannot tell how strong the analogy is for them. The author developed the analogy at length in 2008, in ``The Second Law of Thermodynamics, and Engines of Cognition.''

The post answers only one reading of the religious retort. ``You can't justify your belief in science'' can be a charge of hypocrisy, which the post answers. It can also be a challenge to the standard itself, since science relies on induction, and Hume argued that induction cannot be justified without assuming it. The post does not take up that reading. Our notes on ``The Fallacy of Gray'' discuss that reading and the author's later concession on it.

As in earlier posts, Traditional Rationality is described without naming anyone who holds it, and its falsification rules are Popper's. The research cited for why it is phrased socially is given without the dispute over it: the cheater-detection finding is standard, but whether the social and abstract versions of the task are ``isomorphic'' is disputed, and the finding says nothing about how traditions spread.

In short: a correct answer to ``you do it too,'' based on an analogy to physical law that fits detailed beliefs best, and aimed at a religious retort read only as a charge of hypocrisy.
\end{afterword}
